\documentclass[12pt]{report}

\usepackage{cmap}
\usepackage[T1,T2A]{fontenc}
\usepackage[utf8]{inputenc}
\usepackage[english, russian]{babel}
\usepackage{amssymb}
\usepackage{amsmath}
\usepackage{amsthm}
\usepackage{dsfont}
\usepackage{bm}
\usepackage{diagbox}
\usepackage[left=20mm,right=10mm,top=20mm,bottom=20mm,bindingoffset=2mm]{geometry}
\usepackage{indentfirst}
\usepackage[utf8]{inputenc}
\usepackage{float}
\usepackage[hidelinks]{hyperref}
\usepackage{graphicx}
\usepackage{xcolor}
\usepackage{listings}
\usepackage{minted}
\usepackage{booktabs}
\usepackage{array}
\usepackage{longtable}

\DeclareMathOperator{\N}{\mathbb{N}}
\DeclareMathOperator{\R}{\mathbb{R}}
\DeclareMathOperator{\Z}{\mathbb{Z}}
\DeclareMathOperator{\CC}{\mathbb{C}}
\DeclareMathOperator{\PP}{\mathrm{P}}
\DeclareMathOperator{\Expec}{\mathrm{E}}
\DeclareMathOperator{\Var}{\mathrm{Var}}
\DeclareMathOperator{\Cov}{\mathrm{Cov}}
\DeclareMathOperator{\asConv}{\xrightarrow{a.s.}}
\DeclareMathOperator{\LpConv}{\xrightarrow{Lp}}
\DeclareMathOperator{\pConv}{\xrightarrow{p}}
\DeclareMathOperator{\dConv}{\xrightarrow{d}}

\hypersetup{
	colorlinks=true,
	linkcolor=blue,
	citecolor=blue,
	urlcolor=blue
}

\lstset{language=Python, extendedchars=\true}

\lstdefinestyle{pythonstyle}{
	language=Python,
	backgroundcolor=\color{lightgray},
	commentstyle=\color{green},
	keywordstyle=\color{blue},
	stringstyle=\color{red},
	basicstyle=\ttfamily,
	frame=single,
	breaklines=true
}

\addto\captionsrussian{\renewcommand{\refname}{Список использованных источников}}

\begin{document}
	
	\begin{titlepage}
		\begin{center}
			\large{Федеральное государственное автономное образовательное учреждение высшего образования <<Национальный исследовательский университет ИТМО>>}
		\end{center}
		
		\vspace{15em}
		
		\begin{center}
			\huge{\textbf{Лабораторная работа №5}} \\
			\large{По дисциплине <<Вычислительная математика>>} \\
			\large{Вариант №12}
		\end{center}
		
		\vspace{5em}
		
		\begin{flushright}
			\textit{\large{Выполнил:}} \\
			\large{Студент группы P3218} \\
			\large{Михайлов Дмитрий} \\
			\large{Андреевич} \\
			\textit{\large{Преподаватель:}} \\
			\large{Малышева Татьяна} \\
			\large{Алексеевна}
		\end{flushright}
		
		\vspace{2cm}
		
		\begin{figure}[h]
			\centering
			\includegraphics[width=0.5\linewidth]{image.png}
		\end{figure}
		
		\begin{center}
			Санкт-Петербург \\
			2026 год
		\end{center}
	\end{titlepage}
	
	\tableofcontents
	\newpage
	
	\addcontentsline{toc}{section}{Цель работы}
	\section*{Цель работы}
	Цель лабораторной работы: решить задачу интерполяции и найти значения функции при заданных значениях аргумента, отличных от узловых точек.
	
	\addcontentsline{toc}{section}{Ход работы}
	\section*{Ход работы}
	
	\subsection*{Постановка задачи}
	
	Для варианта №12 из таблицы 1.2 заданы следующие узлы интерполяции:
	\[
	\begin{array}{c|ccccccc}
		x_i & 0{,}50 & 0{,}55 & 0{,}60 & 0{,}65 & 0{,}70 & 0{,}75 & 0{,}80 \\
		\hline
		y_i & 1{,}5320 & 2{,}5356 & 3{,}5406 & 4{,}5462 & 5{,}5504 & 6{,}5559 & 7{,}5594
	\end{array}
	\]
	
	Требуется:
	\begin{itemize}
		\item построить таблицу конечных разностей;
		\item вычислить значение функции в точке \(X_1 = 0{,}523\) по формуле Ньютона;
		\item вычислить значение функции в точке \(X_2 = 0{,}639\) по формуле Гаусса.
	\end{itemize}
	
	Так как узлы являются равноотстоящими, шаг интерполирования равен:
	\[
	h = 0{,}55 - 0{,}50 = 0{,}05.
	\]
	
	\subsection*{Таблица конечных разностей}
	
	По исходным данным построим таблицу конечных разностей.
	
	\[
	\begin{array}{c|c|c|c|c|c|c|c|c}
		i & x_i & y_i & \Delta y_i & \Delta^2 y_i & \Delta^3 y_i & \Delta^4 y_i & \Delta^5 y_i & \Delta^6 y_i \\
		\hline
		0 & 0{,}50 & 1{,}5320 & 1{,}0036 & 0{,}0014 & -0{,}0008 & -0{,}0012 & 0{,}0059 & -0{,}0166 \\
		1 & 0{,}55 & 2{,}5356 & 1{,}0050 & 0{,}0006 & -0{,}0020 & 0{,}0047 & -0{,}0107 & \\
		2 & 0{,}60 & 3{,}5406 & 1{,}0056 & -0{,}0014 & 0{,}0027 & -0{,}0060 & & \\
		3 & 0{,}65 & 4{,}5462 & 1{,}0042 & 0{,}0013 & -0{,}0033 & & & \\
		4 & 0{,}70 & 5{,}5504 & 1{,}0055 & -0{,}0020 & & & & \\
		5 & 0{,}75 & 6{,}5559 & 1{,}0035 & & & & & \\
		6 & 0{,}80 & 7{,}5594 & & & & & &
	\end{array}
	\]
	
	\subsection*{Вычисление значения в точке \(X_1 = 0{,}523\) по первой формуле Ньютона}
	
	Так как точка \(X_1 = 0{,}523\) расположена ближе к началу таблицы, используем первую интерполяционную формулу Ньютона:
	\[
	P_n(x) = y_0 + t \Delta y_0 + \frac{t(t-1)}{2!}\Delta^2 y_0
	+ \frac{t(t-1)(t-2)}{3!}\Delta^3 y_0
	+ \frac{t(t-1)(t-2)(t-3)}{4!}\Delta^4 y_0 + \dots
	\]
	
	Вычислим параметр \(t\):
	\[
	t = \frac{x - x_0}{h} = \frac{0{,}523 - 0{,}50}{0{,}05} = 0{,}46.
	\]
	
	Подставим значения из таблицы конечных разностей:
	\[
	\begin{aligned}
		P(0{,}523) =\;& 1{,}5320
		+ 0{,}46 \cdot 1{,}0036
		+ \frac{0{,}46(0{,}46-1)}{2!}\cdot 0{,}0014 \\
		&+ \frac{0{,}46(0{,}46-1)(0{,}46-2)}{3!}\cdot (-0{,}0008) \\
		&+ \frac{0{,}46(0{,}46-1)(0{,}46-2)(0{,}46-3)}{4!}\cdot (-0{,}0012) \\
		&+ \frac{0{,}46(0{,}46-1)(0{,}46-2)(0{,}46-3)(0{,}46-4)}{5!}\cdot 0{,}0059 \\
		&+ \frac{0{,}46(0{,}46-1)(0{,}46-2)(0{,}46-3)(0{,}46-4)(0{,}46-5)}{6!}\cdot (-0{,}0166).
	\end{aligned}
	\]
	
	Промежуточные слагаемые:
	\[
	\begin{aligned}
		&0{,}46 \cdot 1{,}0036 = 0{,}461656,\\
		&\frac{0{,}46(0{,}46-1)}{2}\cdot 0{,}0014 = -0{,}00017388,\\
		&\frac{0{,}46(0{,}46-1)(0{,}46-2)}{6}\cdot (-0{,}0008) = -0{,}00005100,\\
		&\frac{0{,}46(0{,}46-1)(0{,}46-2)(0{,}46-3)}{24}\cdot (-0{,}0012) = 0{,}00004858,\\
		&\frac{0{,}46(0{,}46-1)(0{,}46-2)(0{,}46-3)(0{,}46-4)}{120}\cdot 0{,}0059 = 0{,}00016911,\\
		&\frac{0{,}46(0{,}46-1)(0{,}46-2)(0{,}46-3)(0{,}46-4)(0{,}46-5)}{720}\cdot (-0{,}0166) = 0{,}00036003.
	\end{aligned}
	\]
	
	Тогда:
	\[
	P(0{,}523) \approx 1{,}994008843.
	\]
	
	Следовательно,
	\[
	\boxed{f(0{,}523) \approx 1{,}9940.}
	\]
	
	\subsection*{Вычисление значения в точке \(X_2 = 0{,}639\) по формуле Гаусса}
	
	Точка \(X_2 = 0{,}639\) находится близко к центральному узлу \(x_0 = 0{,}65\) и левее него, поэтому удобно использовать вторую формулу Гаусса.
	
	Положим:
	\[
	x_0 = 0{,}65,\qquad y_0 = 4{,}5462,\qquad
	t = \frac{x - x_0}{h} = \frac{0{,}639 - 0{,}65}{0{,}05} = -0{,}22.
	\]
	
	Вторая формула Гаусса имеет вид:
	\[
	\begin{aligned}
		P(x)=\;& y_0 + t\Delta y_{-1}
		+ \frac{t(t+1)}{2!}\Delta^2 y_{-1}
		+ \frac{(t-1)t(t+1)}{3!}\Delta^3 y_{-2} \\
		&+ \frac{(t-1)t(t+1)(t+2)}{4!}\Delta^4 y_{-2}
		+ \frac{(t-2)(t-1)t(t+1)(t+2)}{5!}\Delta^5 y_{-3} \\
		&+ \frac{(t-2)(t-1)t(t+1)(t+2)(t+3)}{6!}\Delta^6 y_{-3}.
	\end{aligned}
	\]
	
	Из таблицы конечных разностей:
	\[
	\Delta y_{-1} = 1{,}0056,\quad
	\Delta^2 y_{-1} = -0{,}0014,\quad
	\Delta^3 y_{-2} = -0{,}0020,
	\]
	\[
	\Delta^4 y_{-2} = 0{,}0047,\quad
	\Delta^5 y_{-3} = 0{,}0059,\quad
	\Delta^6 y_{-3} = -0{,}0166.
	\]
	
	Подставим значения:
	\[
	\begin{aligned}
		P(0{,}639)=\;&4{,}5462
		+ (-0{,}22)\cdot1{,}0056
		+ \frac{(-0{,}22)(0{,}78)}{2}\cdot(-0{,}0014) \\
		&+ \frac{(-1{,}22)(-0{,}22)(0{,}78)}{6}\cdot(-0{,}0020) \\
		&+ \frac{(-1{,}22)(-0{,}22)(0{,}78)(1{,}78)}{24}\cdot0{,}0047 \\
		&+ \frac{(-2{,}22)(-1{,}22)(-0{,}22)(0{,}78)(1{,}78)}{120}\cdot0{,}0059 \\
		&+ \frac{(-2{,}22)(-1{,}22)(-0{,}22)(0{,}78)(1{,}78)(2{,}78)}{720}\cdot(-0{,}0166).
	\end{aligned}
	\]
	
	После вычислений получаем:
	\[
	P(0{,}639) \approx 4{,}325103662.
	\]
	
	Следовательно,
	\[
	\boxed{f(0{,}639) \approx 4{,}3251.}
	\]
	
	\subsection*{Итоги вычислительной части}
	
	В ходе вычислительной части была построена таблица конечных разностей и найдены приближённые значения функции в заданных точках:
	\[
	\boxed{f(0{,}523) \approx 1{,}9940}
	\]
	по первой формуле Ньютона и
	\[
	\boxed{f(0{,}639) \approx 4{,}3251}
	\]
	по второй формуле Гаусса.
	
	\newpage
	
	\addcontentsline{toc}{section}{Блок-схемы используемых методов}
	\section*{Блок-схемы используемых методов}	
	
	\begin{figure}[H]
		\centering
		\includegraphics[width=0.3\linewidth]{Рисунок1.png}
		\caption{Блок-схема многочлена Лагранжа}
	\end{figure}
	
	\begin{figure}[H]
		\centering
		\includegraphics[width=0.3\linewidth]{Рисунок2.png}
		\caption{Блок-схема многочлена Ньютона}
	\end{figure}
	
	
	\addcontentsline{toc}{section}{Листинг программы}
	\section*{Листинг программы}
	\href{https://github.com/mysticslippers/math_comp_archive/tree/main/Lab5_comp_math/code}{Ссылка на репозиторий с кодом.}
	
	\addcontentsline{toc}{section}{Результат выполнения программы}
	\section*{Результат выполнения программы}
	
	
	\begin{figure}[H]
		\centering
		\includegraphics[width=0.6\linewidth]{screen1.png}
		\caption{Пример выполнения программы.}
	\end{figure}
	
	\begin{figure}[H]
		\centering
		\includegraphics[width=0.6\linewidth]{screen2.png}
		\caption{Пример выполнения программы.}
	\end{figure}
	
	\begin{figure}[H]
		\centering
		\includegraphics[width=0.6\linewidth]{screen3.png}
		\caption{Пример выполнения программы.}
	\end{figure}
	
	\begin{figure}[H]
		\centering
		\includegraphics[width=0.6\linewidth]{screen4.png}
		\caption{Пример выполнения программы.}
	\end{figure}
	
	\begin{figure}[H]
		\centering
		\includegraphics[width=0.6\linewidth]{screen5.png}
		\caption{Пример выполнения программы.}
	\end{figure}
	
	\begin{figure}[H]
		\centering
		\includegraphics[width=0.8\linewidth]{graph.png}
		\caption{Вывод графиков функций.}
	\end{figure}
	
	
	\addcontentsline{toc}{section}{Вывод}
	\section*{Вывод}
	
	\noindent
	В ходе выполнения лабораторной работы была решена задача интерполяции функции по таблично заданным данным. Для варианта №12 была построена таблица конечных разностей, после чего вычислены приближённые значения функции в заданных точках с использованием интерполяционных формул Ньютона и Гаусса.
	
	\medskip
	
	\noindent
	В результате вычислений получены следующие значения:
	\[
	f(0{,}523) \approx 1{,}9940,
	\qquad
	f(0{,}639) \approx 4{,}3251.
	\]
	
	\medskip
	
	\noindent
	Также в программной части были реализованы методы Лагранжа, Ньютона с разделёнными разностями и Гаусса. Программа позволяет задавать исходные данные несколькими способами, формировать таблицу конечных разностей, вычислять значение функции в произвольной точке и строить графики интерполяционных многочленов.
	
	\medskip
	
	\noindent
	По результатам работы можно сделать вывод, что интерполяционные методы позволяют эффективно находить приближённые значения функции внутри заданного интервала по известным табличным значениям. Для равноотстоящих узлов удобно применять формулы Ньютона и Гаусса, так как они используют конечные разности и дают точные и наглядные результаты.
	
	
\end{document}
